%%%%%%%%%%%%%%%%%%%%%
% Documento maestro %
%%%%%%%%%%%%%%%%%%%%%
\documentclass{fime}

%%%%%%%%%%%%%%%%%%%%%%%%%%%%%%%%%%%%%%%%%%%
% Cargando paquetes y definiendo opciones %
%%%%%%%%%%%%%%%%%%%%%%%%%%%%%%%%%%%%%%%%%%%
% Aquí puedes cargar los paquetes que vas a usar. La clase
% fime ya incluye babel, inputenc, graphicx y los de la AMS.
% Cargar un paquete está a tu libertad (y responsabilidad).
\usepackage{hyperref}
    \hypersetup{breaklinks=true,colorlinks=true,
        linkcolor=black,citecolor=black,urlcolor=black}
\usepackage{datetime} % monthname
\usepackage{todonotes} % temporaly use for todo notes
\usepackage{algorithm}
\usepackage{algpseudocode}

%%%%%%%%%%%%%%%%%%%%%
% Definiendo campos %
%%%%%%%%%%%%%%%%%%%%%
\def\titulo{Optimal Location of Car Wreck Adjusters}
\def\autor{Luis Alberto Maltos Ortega}
\def\matricula{1390200}
\def\grado{Maestría en Ciencias en Ingenie\'ia de Sistemas}
% En caso de que el grado tenga orientación o especialidad
% llenar el siguiente campo dejando un ESPACIO INICIAL,
% en caso contrario, dejar vacío
\def\orientacion{ }
\def\fecha{\monthname,  \the\year} % Coloca el mes con mayúscula inicial

\def\asesor{Dr. Roger Z. R\'ios Mercado}
\def\revisorA{Mar\'ia Ang\'elica Salazar Aguilar}
\def\revisorB{Mar\'ia Guadalupe Villareal Marroqu\'in}
% En el caso de que tu tesis sea de doctorado
% activa la variable cambiándola a \doctoradotrue
% y define tus otros dos revisores
\newif\ifdoctorado\doctoradofalse
\def\revisorC{Nombre del revisor C}
\def\revisorD{Nombre del revisor D}
% El visto bueno siempre va
\def\vobo{Dr. Simón Martínez Martínez}

\includeonly{
  Misc/Cover,
  Misc/Dedicatory,
  Misc/Acknowledgements,
  Misc/Abstract,
  Introduction/Main,
  RelatedWork/Main,
  ProblemStatement/Main,
  Methodology/Main,
  Models/Main,
  HeuristicProcedures/Main,
  Evaluation/Main,
  Conclusions/Main,
  Appendices/Jarvis
}
%%%%%%%%%%%%%%%%%%%%%%%
% Inicia el documento %
%%%%%%%%%%%%%%%%%%%%%%%
\begin{document}

\frontmatter
\pagestyle{main}

%%% Incluye PortillasM si tu tesis es de Maestría
%%% y PortillasD si es de doctorado.
% Portadas (Maestría)

\def\uanl{Universidad Autónoma de Nuevo León}
\def\fime{Facultad de Ingeniería Mecánica y Eléctrica}
\def\depg{Subdirección de Estudios de Posgrado}
\def\snnl{San Nicolás de los Garza, Nuevo León}

%%%%%%%%%%%%%%%%%%%%%%%%
% Primer portada: UANL %
%%%%%%%%%%%%%%%%%%%%%%%%
\thispagestyle{empty}

\begin{scshape}
\begin{center}
	{\Large \uanl} \\[5mm]
	{\large \fime} \\[5mm]
	{\large \depg}
	\vskip 15mm
	\includegraphics[height=55mm]{uanl}
	\vskip 12mm
	\begin{tabular}{p{11cm}}
		\centering
		{\large \titulo}
	\end{tabular}
	\vskip 7mm
	{por}\\[7mm]
	{\large \autor}\\[7mm]
	{como requisito parcial para obtener el grado de}\\[3mm]
	\MakeUppercase{\grado}\\
	\orientacion
	\vfill
	\fecha
\end{center}
\end{scshape}

%%%%%%%%%%%%%%%%%%%%%%%%%
% Segunda portada: FIME %
%%%%%%%%%%%%%%%%%%%%%%%%%
\newpage
\thispagestyle{empty}

\begin{scshape}
\begin{center}
	{\Large \uanl} \\[5mm]
	{\large \fime} \\[5mm]
	{\large \depg}
	\vskip 16mm
	\includegraphics[height=50mm]{fime}
	\vskip 16mm
	\begin{tabular}{p{11cm}}
		\centering
		{\large \titulo}
	\end{tabular}
	\vskip 7mm
	{por}\\[7mm]
	{\large \autor}\\[7mm]
	{como requisito parcial para obtener el grado de}\\[3mm]
	\MakeUppercase{\grado}\\
	\orientacion
	\vfill
	\fecha
\end{center}
\end{scshape}

%%%%%%%%%%%%%%%%%%%%%%%%%%%%%
% Carta del comité de tesis %
%%%%%%%%%%%%%%%%%%%%%%%%%%%%%
\newpage
\thispagestyle{empty}
\enlargethispage{5mm}

\begin{center}
{\bf \large \uanl} \\
{\bf \fime} \\
{\bf \depg}
\end{center}
\vskip 4mm

Los miembros del Comité de Tesis recomendamos que la Tesis <<\titulo>>, realizada por el alumno \autor, con número de matrícula \matricula, sea aceptada para su defensa como requisito parcial para obtener el grado de \grado\orientacion.
\ifdoctorado{\vskip 10mm}\else{\vskip 8mm}\fi

\begin{center}
El Comité de Tesis\\
\ifdoctorado{\vskip 15mm}\else{\vskip 25mm}\fi

\ifdoctorado{%%%
\begin{tabular}{p{37mm}p{21mm}p{12mm}p{21mm}p{37mm}}
	\cline{2-4}
	& \multicolumn{3}{c}{\asesor} & \\
	& \multicolumn{3}{c}{Asesor}  & \\[15mm]
	\cline{1-2} \cline{4-5}
	\multicolumn{2}{c}{\revisorA} & & \multicolumn{2}{c}{\revisorB} \\
	\multicolumn{2}{c}{Revisor}   & & \multicolumn{2}{c}{Revisor}   \\[17mm]
	\cline{1-2} \cline{4-5}
	\multicolumn{2}{c}{\revisorC} & & \multicolumn{2}{c}{\revisorD} \\
	\multicolumn{2}{c}{Revisor}   & & \multicolumn{2}{c}{Revisor}   \\[2mm]
	& \multicolumn{3}{c}{Vo. Bo.} & \\[14mm]
	\cline{2-4}
	& \multicolumn{3}{c}{\vobo}   & \\
	& \multicolumn{3}{c}{\depg}   & \\ &&&&
\end{tabular}
}\else{%%%
\begin{tabular}{p{37mm}p{21mm}p{12mm}p{21mm}p{37mm}}
	\cline{2-4}
	& \multicolumn{3}{c}{\asesor} & \\
	& \multicolumn{3}{c}{Director}  & \\[19mm]
	\cline{1-2} \cline{4-5}
	\multicolumn{2}{c}{\revisorA} & & \multicolumn{2}{c}{\revisorB} \\
	\multicolumn{2}{c}{Revisor}   & & \multicolumn{2}{c}{Revisor}   \\[2mm]
	& \multicolumn{3}{c}{Vo. Bo.} & \\[17mm]
	\cline{2-4}
	& \multicolumn{3}{c}{\vobo}   & \\
	& \multicolumn{3}{c}{\depg}   & \\ &&&&
\end{tabular}
}\fi%%%

\vfill

\snnl, \MakeLowercase{\fecha}

\end{center}

% Dedicatoria

\thispagestyle{empty}
\vspace*{17mm}

\begin{flushright}
\begin{itshape}

Aquí puedes poner tu dedicatoria\\
si es que tienes una.\bigskip\bigskip

Si no tienes una, puedes borrar\\
la línea \verb+% Dedicatoria

\thispagestyle{empty}
\vspace*{17mm}

\begin{flushright}
\begin{itshape}

Aquí puedes poner tu dedicatoria\\
si es que tienes una.\bigskip\bigskip

Si no tienes una, puedes borrar\\
la línea \verb+% Dedicatoria

\thispagestyle{empty}
\vspace*{17mm}

\begin{flushright}
\begin{itshape}

Aquí puedes poner tu dedicatoria\\
si es que tienes una.\bigskip\bigskip

Si no tienes una, puedes borrar\\
la línea \verb+\include{Dedicatoria}+ en el\\
archivo \texttt{MiTesis.tex} pues no es obligatoria.

\end{itshape}
\end{flushright}

+ en el\\
archivo \texttt{MiTesis.tex} pues no es obligatoria.

\end{itshape}
\end{flushright}

+ en el\\
archivo \texttt{MiTesis.tex} pues no es obligatoria.

\end{itshape}
\end{flushright}



\tableofcontents
\listoffigures
\listoftables
\listoftodos
\todo{Remove list of to do's when complete}

%Acknowledgements

\chapter{Acknowledgments}
\markboth{Acknowledgments}{}

%Aquí puedes poner tus agradecimientos.
%(No olvides agradecer a tu comité de tesis,
I want to especially thank
Prof. Roger Z. Rios
For giving me the opportunity to work with him
and be my advisor,
for guiding me in my academic development
and for the effort and support he gave me
I also would like to thanks him
for his patience,
without his support
this work would not have been possible. 

Thanks to my thesis committee
Prof. Ma. Angelica Salazar Aguilar
and Prof. Ma. Guadalupe Villareal Marroquin,
for their comments and suggestions
that made throughou his work,
and for help me
to complete this work,

%a tus profesores,
to all professors of the PISIS
for their courses so complete and enriching,

%a la facultad
to the FIME
for the facilities granted
and ease the paperwork,

%y a CONACyT
%en caso de que hallas sido beneficiado con una beca).
to Consejo Nacional de Cíencia y Tecnología (CONACyT)
for the Graduate Fellowship
and the research project grant 2011-1-166397
with which it was possible the realization of my studies
and attendance at national conferences
to present part of this work.

%Resumen

\chapter{Abstract}
\markboth{Abstract}{}

{\setlength{\leftskip}{10mm}
\setlength{\parindent}{-10mm}

\autor.

Candidato para obtener el grado de \grado\orientacion.

\uanl.

\fime.

Título del estudio: \textsc{\titulo}.

\noindent Número de páginas: \pageref*{lastpage}.}

%%% Comienza a llenar aquí
\paragraph{Problem Description:}
When a traffic accident occurs
in cities with a large traffic flow
the roads surrounding the crash site
are affected by traffic congestions.
In some countries, such as Mexico,
even small accidents are troublesome
due to the fact
that a claim adjuster
from the car insurance company
must arrive to the site
and document the accident
before the vehicle may be moved
as required by law.
Thus,
in this particular setting,
the adjusters location
becomes a key factor in providing timely service.

\paragraph{Objectives and method of study:}
%Aquí debes poner tus objetivos y métodos de estudio.
The aim of this thesis is to
\begin{itemize}
\item Provide quantitative tools
with scientific support
for the optimal location
of car wreck adjusters.
\item Develop
adequate mathematical models
for representing
some of the important company concerns.
Design and develop efficient solution techniques
for handling real-world instances of the problem.
\item Assess
the performance of the proposed techniques
based on an appropriate experimental design.
\end{itemize}

\paragraph{Contributions and Conclusions:}
%Y aquí tus contribuciones y conclusiones.
%(También es parte del formato).
The most important contributions of this work
are the following:
\begin{itemize}
\item Two integer programming models
  are introduced
  with the aim of
  minimizing average response time.
\item The second model
  takes into account
  special features
  to reduce the number of variables.
\item A scatter search heuristic
  was designed, implemented, and tested
  on a wide set of instances
  with very good results.
\end{itemize}

\bigskip\noindent\begin{tabular}{lc}
\vspace*{-2mm}\hspace*{-2mm}Firma del asesor: & \\
\cline{2-2} & \hspace*{1em}\asesor\hspace*{1em}
\end{tabular}




\mainmatter
\pagestyle{fime}

%%% Haz un documento para cada capítulo
\chapter{Introduction}
The main idea behind this thesis
is to develop models and methods
for improving the service offered by car insurance agents,
helping them arrive to accident sites sooner,
and determining the number of adjusters required
to perform the service within the desired standards.

\section{Problem Statement}
A car insurance company
in countries such as Mexico
faces the daily issue
of how to locate
their insurance agents (adjusters),
in such a way
that they provide
the best possible service
to their customers.
Typically
the quality of service
is measured
by how quickly
the adjusters arrive
to the places
where an incident has occurred.

This problem
can be seen
as an emergency location system, % add cite
in which
it is desired to ``cover''
the territory under study
(city intersections)
by a fixed number of adjusters
so as to reduce
the average time
of arrival to accidents.
Another issue
that makes the problem
even more complex
is the fact
that accidents occur randomly
and decisions
on where to place the adjusters
must be made before
accidents occur.

\section{Background}
This is a practical problem that has not been studied,
ya que este sistema de emergencia,
no se presenta en muchos paises,
sin embargo, hay sistemas similares,
tal es el caso de ambulancias, bomberos, patrullas,
para los cuales existen una gran cantidad de articulos al respecto,
pero las caracteristicas tratadas en ellos,
difieren de nuestro problema practico,
por los siguientes factores
tipo de espacio,
radio cobertura,
multiperiodo,
tiempo de servicio,
tipo de llegadas

Berman propone un problema de localizacion para sistemas de emergencia,
y lo llama The Queue Stochastic p-Median Problem.

\section{Motivation}
When a car accident occurs,
traffic congestion starts to pile up.
This is because
customers are not allowed
to move their vehicles until
the adjuster arrives.
The adjuster must record and determine
the causes of the accident,
in order to
move the car from the accident area
and restore the flow.

\section{Objectives}
The aim of this work
is to support decision-making,
regardin location of insure agents
to attend car wrecks.

The main goal is
to determine the optimal bases (locations)
for placing the insurance company adjusters,
so as to minimize
the average or maximum response time
from customer calls
when accidents occur.

\section{Organization}
This thesis is organized as follows:
Chapter \ref{ch:RelatedWork}
presents an brief literature review
of diverse approaches of Emergency Service Systems (ESS),
starting with the simple deterministic models
and ending with the hypercube model and simulation models.
The framework is discussed in Chapter \ref{ch:Framework}
the problem statement and two models are presented.
Chapter \ref{ch:Heuristic} describes the heuristic method proposed
and their components.
Chapter \ref{ch:Experiments} contains
the experiments made to the models,
and to the heuristics to tune them.
Conclusions, contributions, 
and directions for future research
are highlighted in Chapter \ref{ch:Conclusions}.
Additionally, Appendix \ref{ch:Jarvis}
contains a summary of how to approximate the hypercube results,
and together with Appendix \ref{ch:Berman}
completes the explanation of how to estimate
the response times of a solution.
%cedure discussed in Chapter 5
%Chapter 6 describes
%Chapter 7 contains a
%Conclusions, contributions, 
%and directions for future research are highlighted in Chapter 8
%Additionally, Appendix A


\chapter{Related Work}
In this chapter
the first works
related to location of vehicles
are introduced.
% source: brotcorner2003ambulance
The models are classified in two main categories,
deterministic and probabilistic.
Deterministic models
are used at the planning stage
and ignore stochastic considerations
regarding the availability of vehicles.
Probabilistic models
reflect the fact that
vehicles operate as servers in a queueing system
and cannot always answer a call.
% eos

\section{Covering models}
% source: brotcorner2003ambulance
the location set covering model (LSCM)
introduced Toregas et al. \cite{toregas1971location}
\todo[size=\tiny,caption={Add ref}]{include toregas1971location in refs}
the aim
is to minimize
the number of ambulances needed
to cover all demand points.
This model
ignores several aspects of real-life problems,
the most important
probably being that
once an ambulance is dispatched,
some demand points
are no longer covered.
However,
provide
a lower bound
on the number of ambulances
required to ensure full coverage.

In the maximal covering location problem (MCLP)
originally proposed by Church and ReVelle \cite{church1974maximal}
\todo[size=\tiny,caption={Add ref}]{include church1974maximal in refs}
is an alternative approach
proposed to counter some of the shortcomings of the LSCM
is to maximize population coverage
subject to limited ambulance availability.
% eos

\subsection{P-Median Models}
% source: golberg1990validating
The first approaches
incorporate median type objectives,
the aim of this problem is
\todo[inline]{Explicar en que consisten estos modelos}
Examples of such models are
the deterministic p-median of Hakimi \cite{hakimi1964optimum},
the stochastic p-median and the vector assignment p-median of Church and Weaver
\cite{weaver1983computational,weaver1985median}
and the capacitated p-median of Schilling and Pirkul \cite{pirkul1988siting}

The strength of these models
is that optimal solution procedures
have been developed
and can accommodate practical problems.
Weaknesses include assumptions
such as noncooperation between vehicles,
the probability of each system state is known
and the fraction of call
served by the closest, second closest, etc.
is known for each zone.
% eos

\section{Hypercube Models}
The hypercube model
was developed by Larson \cite{larson1974hypercube}
%and extended by others (see Swersey \cite{swersey1994deployment})
% source: galvao2008emergency
is not an optimization model; 
it is only a descriptive model
that permits the analysis of scenarios.

Given a system configuration,
the hypercube model
is able to evaluate
a variety of performance measures relevant for decision-making,
either region-wide or for each server or region. 
These include
server workloads,
mean user response times,
fraction of dispatches of each server to each region,
among others.
\todo[inline]{Include definitions of the model}

\subsection{Calibration Process}
% source: galvao2008emergency
% Calibration process of the mean service times
In certain EMSs
and other emergency systems,
travel times may represent
a considerable part of service times.
In such cases,
it may be advisable
to adjust the service times
by means of a calibration process,
which can be performed
using a simple iterative procedure
that is proposed by Berman et al. \cite{berman1987stochastic}.
Basically,
the procedure consists of
verifying
if there are
significant differences among the input mean service times
and the output mean service times (computed by the hypercube model).
In this case,
the hypercube is solved
using the computed mean service times as inputs,
until
the differences
among input and output values
are sufficiently small.
This procedure is called
a calibration process.
Note that
it takes into account
that the mean travel time
depends on the location of the user
and the identity of the server.
Empirical experiments show
that this procedure
usually converges in two or three iterations,
for a reasonably accurate estimation of the mean service times,
although a formal proof of the convergence of the method
is apparently not available in the literature.

\section{Probabilistic location models for planning ESSs}
% source: galvao2008emergency
Probabilistic location models
are optimization models
that permit a more accurate planning
of EMSs at the strategic level.
They were initially developed
based on simplifying assumptions of independence
and non-cooperation among servers,
which are seldom encountered in practice.
% eos
% source: galvao2008emergency
The location of facilities (servers) in ESSs
has as its main objective
the provision of coverage to demand areas.
The definition of
probabilistic location models for planning these systems
was a natural extension
of their deterministic equivalents,
the location models with covering constraints.
The notion of coverage
implies
the definition of a service distance (time),
which is the critical distance (time)
beyond which a demand area is considered not covered.
A demand area is
therefore
considered covered
if
it is within a predefined critical distance (say D)
from at least one of the existing facilities.
%The simplest deterministic covering model
%corresponds to the Location Set Covering Problem (LSCP),
%which provides coverage for every demand area under consideration.
%The provision of total coverage, however,
%may prove to be economically infeasible,
%in the sense that the number of servers required
%may not be compatible with the resources available to the decision-maker.
%The Maximal Covering Location Problem (MCLP)
%was defined by Church and ReVelle (1974) within this context.
%In this case the objective is to locate a number of facilities
%(say p facilities) that is compatible with the resources available,
%such that the maximal possible population of a given geographical region
% is covered within D.

A limitation of the deterministic models
is that
they assume
that servers are available when requested,
which is not always true in practical situations.
Congestion in emergency services,
which may cause
the unavailability
of servers located
within the critical distance
when a call is placed,
lead to
the development of
a second generation of location covering models
focused on additional coverage.
These models
emphasize the importance of additional coverage's
for the demand areas,
given the possibility that
in congested systems
the first server,
possibly
the only server in a particular coverage area,
might not be available when requested.
Several such models were developed,
as for example in Daskin and Stern (1981),
Eaton et al. (1981),
Hogan and ReVelle (1986)
and Batta and Mannur (1990).
Probabilistic covering models
are a natural extension
of the second-generation models.
The initial problems
defined within this context
were
the Maximum Expected Covering Location Problem (MEXCLP)
and
the Maximum Availability Location Problem (MALP).
%Extensions of these problems
%lead to the use
%of the hypercube queuing model
%in corresponding solution methodologies.
% eos

\section{Simulation Models}
% source: golberg1990validating
Simulation models
can be formulated with a great deal of detail
and hence can be validated.
% eos

% source: goldberg1990simulation
Simulation
is used for evaluating EMS system performance in numerous papers
as in Savas \cite{savas1969simulation},
Berlin and Liebman \cite{berlin1974mathematical},
and Swoveland et al. \cite{swoveland1973simulation},

Their provides in general
a wealth of output measures,
drawback is that
they are rarely used
because of high runtime and data collection costs.
These models have some objectionable assumptions,
but successful applications do exist in the literature.

\todo[inline]{Include
-Ambulance location
-Cooperative service
-Stocatic Queue p-Median model}

\chapter{Problem Statement}
In this chapter we describe the problem
mentioned previously

% Ideal Situation
The location and dispatching policies
used by insurance agencies
should aim
to arrive to incidents
early as possible,
so that
the workloads for adjusters are balanced.
The location policies
should be optimized for service time
but also should consider
cooperatively
based on workloads.

% Problem
However,
current insurance agency policies
are empirical,
and does not consider cooperative.
Taking more time to arrive at incidents,
making the insurance agency
less competitive,
and generating more traffic congestion.

% Financial Costs
% Evidence
% Solution
%  Benefits
% Summary
% Thesis statement
% Conceptual Problems

\chapter{Methodology}
In this chapter
we describe the solution methodology proposed
to solve the described problem.

A first model was created
base on proposed by Goldberg \cite{goldberg1990validating}
for which,
we made some relaxations
instead obtaining a linear model.

Proposed a second model
that performs an additional simplification,
considering that
it is unlikely
assign an adjuster
to a demand point
covered previously
by more adjusters,
omitting allocation variables
and adding restrictions
to guarantee
the correct order of allocation.

Because the problem stochasticity
can not evaluate solutions accurately,
so the Hypercube model is used.

Several functions were evaluated
for constructive algorithm
but because presented cooperativeness
it is difficult
to approximate the final results
in a partial solution
with a greedy function.

To improve these solutions,
a scatter-search was implemented
with a dynamic reference set.
We opted for a dynamic reference set
due the number of solutions generated
in each combination is big.

\chapter{Models}

% Validating and applying a model 
% for locating emergency medical vehicles in Tucson, AZ
% Goldberg 1990

\section{Model A}

This model is based
on the model
proposed by Goldberg \cite{goldberg1990validating},
and it contains assignment variables
for all possible orders.

The following assumptions are considered in the model:
\begin{itemize}
\item The probability that an adjuster is busy
  is $\rho$ and is unaffected by the state of the system
\item There is a strict ordering of the bases preferred for each zone
  that does not depend
  on the current state of the system
\item All calls are answered
  by an adjuster originating from its base,
  not en route back to the base
\item The arrival of calls to the system
  follows a stationary distribution
\item The model is presented
  using a 0-queue assumption.
\end{itemize}

Sets and indexes:
\begin{itemize}
\item \textit{n} number of demand points
\item \textit{m} number of potential sites to locate adjusters/facilities
\item \textit{p} number of available adjusters
\item \textit{i} index for demand points;
  $i \in V = \{1,2,\ldots,n\}$ 
\item \textit{j} index for potential site for adjusters/facilities
  $j \in W = \{1,2,\ldots,m\}$
\item $k$ possible order;
  $k \in K = \{1,2\ldots,p\}$
\item $S_{ij} = \{
  r \in W | \mbox{ site } r \mbox{ is preferred before site } j
  \mbox{ for demand point } i
  \}$
\end{itemize}

Parameters:
\begin{itemize}
\item $\rho$ is the utilization of each adjuster
\item $t_{ij}$ is the expected travel time
  between demand point \textit{i}
  and potential site \textit{j}.
\item $h_{ij}^{k}$ is the probability
  that $j$ serves point $i$
  given that
  is the $k$-th preferred.
\end{itemize}

Variables:
\begin{itemize}
\item $x_j =
  \begin{cases} 
    1 & \mbox{if an adjuster is at potential site } j \\
    0 & \mbox{otherwise}
  \end{cases}$
\item $y_{ij}^{k} =
  \begin{cases} 
    1 & \mbox{if adjuster in } j \mbox{, is the }
    k\mbox{-th to cover demand point } i \\
    0 & \mbox{otherwise}
  \end{cases}$
\end{itemize}

Model
{\small
  \begin{equation}
    \min \, \sum_{j=1}^{m}{
      \sum_{k=1}^{p}{
        \sum_{i=1}^{n}{
          h_{ij}^{k}t_{ij}y_{ij}^{k}
        }
      }
    }
  \end{equation}
}
Minimize the average expected response time
\begin{align}
  \label{eq:2}
  \sum_{j \in W}{x_j}
  & = p
\end{align}
Only locate $p$ adjusters
\begin{align}
  \label{eq:3}
  \sum_{j \in W}{y_{ij}^{k}}
  & = 1
  & i \in V, k
  &\in K
\end{align}
Each demand point $i$ is covered by an adjuster on each order $k$
\begin{align}
  \label{eq:4}
  y_{ij}^{k}
  & \leq x_j
  & i \in V,j \in W, k
  &\in K
\end{align}
Assigned only if there is an adjuster located at $j$
\begin{align}
  \label{eq:5}
  \sum_{k = 1}^{p}{
    y_{ij}^{k}
  }
  & \leq x_j
  & i \in V, j 
  & \in W
\end{align}
Only assign in one order for each adjuster located 
\begin{align}
  y_{ij}^{k} 
  & \leq \sum_{r\in S_{ij}}{y_{ir}^{k-1}}
  & i \in V,j \in W, k
  & \in K/\{1\}
\end{align}
Assign \textit{j} to cover \textit{i} in order $k$
only if
the assign of order $k-1$
was made for some $r \in S_{ij}$
\begin{align}
  x_{j}
  & \in \{0,1\}
  & j 
  & \in W \nonumber
  \\
  y_{ij}^{k}
  & \in \{0,1\}
  &  i \in V,j \in W,k
  &\in K \nonumber
\end{align}

%This model was tested 
%with a benchmark instance of 30 points, 
%which sets $V$ and $W$ are the same.
%For values of $p = 1,\ldots,7$
%the optimal solution was found in less than one minute.
%
%For other benchmark instances
%was impossible to find the solution
%because they have more than 300 points,
%and the proposed model grows to over one million variables.
%  For some large instances ($n = 50$ , $m = 30$)%
\todo[inline]{Add Scope}
\todo[inline]{Computational Results}

% A model for the Stochastic Queue p-Median Problem
% Created by Luis Maltos & Roger Rios
% 2015

\section{Model B}

This model
was developed with the idea
that is unlikely that the farthest adjusters
serve demand points
on cases where the system does not become congested.
In these cases
we can make the assumption that
the probability of being served by the $\ell$-th 
adjuster is almost zero,
where $\ell$ is large enough but less than $p$.

Parameters:
\begin{itemize}
\item $M$ is a large integer
\item $\ell$ the number of allowed adjusters per demand point
\item $a_{ik}$ the $k$-th preferred location server
  regarding the point $i$.
\end{itemize}

Variables:
\begin{itemize}
\item $z_j$ the number of adjusters placed at site $j$
\item $y_{ij}^k$ if adjuster in $j$,
  is the $k$-th to cover demand point $i$
\end{itemize}

The objective,
and the constraints \ref{eq:2}-\ref{eq:5} 
are practically the same as in model A,
with the difference that
binary variables $x_j$ from Model A
are replaced by integer variables $z_j$
inspired by the results of Berman \cite{berman1987stochastic},
and the addition of the following binary variables
\begin{itemize}
\item $u_{ij} = 
  \begin{cases}
    1 & \mbox{if the number of adjusters between } i 
    \mbox{ and } j \mbox{, inclusive, is less than } \ell \\
    0 & \mbox{otherwise}
  \end{cases}$
\item $v_{ij} = 
  \begin{cases} 
    1 & \mbox{if the number of adjusters between } i
    \mbox{ and } j \mbox{, is less than } \ell - 1 \\
    0 & \mbox{otherwise}
  \end{cases}$
\end{itemize}

\begin{align}
  \sum_{r = 1}^{k \mid a_{ik}=j}{
    z_{a_{ir}}
  }
  + (p-\ell) u_{ij}
  & \leq p
  & i \in V,
  & j \in W
\end{align}
Checks if
the number of adjusters between $i$ and $j$,
inclusive,
is less or equal than $\ell$
\begin{align}
  \sum_{r = 1}^{k \mid a_{ik}=j}{
    z_{a_{ir}}
  }
  + M u_{ij}
  & \geq \ell+1
  & i \in V,
  & j \in W
\end{align}
Checks if
the number of adjusters between $i$ and $j$,
inclusive,
is greater than $\ell$
\begin{align}
  \sum_{k = 1}^{\ell}{
    y_{ij}^{k}
  }
  + M (1 - u_{ij})
  & \geq z_j
  & i \in V,
  & j \in W
\end{align}
Assign $z_j$ times $j$ to $i$ if $u_{ij} = 1$
\begin{align}
  \sum_{r = 1}^{k \mid a_{i(k+1)}=j}{
    z_{a_{ir}}
  }
  + (p-(\ell-1)) v_{ij}
  & \leq p
  & i \in V,
  & j \in W
\end{align}
Checks if
the number of adjusters between $i$ and $j$,
is less than $\ell$
\begin{align}
  \sum_{r = 1}^{k \mid a_{i(k+1)}=j}{
    z_{a_{ir}}
  }
  + M v_{ij}
  & \geq \ell
  & i \in V,
  & j \in W
\end{align}
Checks if
the number of adjusters between $i$ and $j$,
is greater or equal than $\ell$
\begin{align}
  \sum_{k=1}^{k}{
    y_{ij}^{k}
  }
  + M (1 - v_{ij} + u_{ij})
  & \geq \ell
  - \sum_{r = 1}^{k \mid a_{i(k+1)}=j}{
    z_{a_{ir}}
  } 
  &  i \in V, j 
  & \in W
\end{align}
Assign $j$ to $i$
the times remaining to complete $\ell$ assignments
\begin{align}
  \sum_{k=1}^{k}{
    y_{ij}^{k}
  }
  - M (1 - v_{ij} + u_{ij})
  & \leq \ell 
  - \sum_{r = 1}^{k \mid a_{i(k+1)}=j}{
    z_{a_{ir}}
  }
  & i \in V,
  & j \in W
\end{align}
Assign $j$ to $i$
the times remaining to complete $\ell$ assignments
\begin{align}
  y_{ij}^{k}
  & \leq u_{ij}
  + v_{ij}
  & i \in V,
  & j \in W
\end{align}
Assign $j$ to $i$
only if
is in the first $\ell$ adjusters near $i$
\begin{align}
  z_j
  & \in \{0,1,\ldots,p\}
  & j
  & \in V \nonumber
  \\
  y_{ij}^{k} 
  & \in \{0,1\}
  & i\in V,j\in W,k
  & \in I \nonumber
  \\
  u_{ij},v_{ij}
  & \in \{0,1\}
  & i \in V,j
  & \in W \nonumber
\end{align}


\chapter{Heuristic Procedures}
Because the problem stochasticity
can not evaluate solutions accurately,
so the Hypercube model is used
in conjunction with
the service mean time calibration process
(see Berman et al. \cite{berman1987stochastic})
to obtain the main parameters.

\section{Proposed Metaheuristic}
The Scatter Search (SS)
is an evolutionary approach
introduced by Glover \cite{glover1977heuristics}
as a metaheuristic for integer programming.
It is based on diversifying the search
through the solution space.
It operates on a set of solutions,
named the Reference Set,
formed by good and diverse solutions of the main population.
These solutions are combined
with the aim of generate
new solutions with better fitness,
while maintaining diversity.
The General Scheme can be seen in Figure \ref{fig:SS_GS}.
\begin{figure}
  \centering
  \begin{tikzpicture} [
      auto,
      cloud/.style = { ellipse,
        draw=blue,
        fill=blue!20,
        node distance=3cm,
        minimum height=2em },
      decision/.style = { diamond,
        draw=blue,
        thick,
        fill=blue!20,
        text width=5em,
        text badly centered,
        inner sep=1pt,
        rounded corners },
      block/.style    = { rectangle,
        draw=blue,
        thick, 
        fill=blue!20,
        text width=10em,
        text centered,
        rounded corners,
        minimum height=2em },
      line/.style     = { draw, thick, ->, shorten >=2pt },
    ]
    % Define nodes in a matrix
    \matrix [column sep=5mm, row sep=10mm] {
      & \node [block] (init) {Diversification Generation};
      & &
      \\
      & \node (null1) {};
      & \node [block] (update) {\textsf{Update Reference Set}};
      & 
      \\
      & \node [cloud] (improve) {\textsf{Improvement}};
      & \node [decision] (newsol) {\textsf{Is there a new solution?}};
      & \node [cloud] (stop) {\textsf{Stop}};
      \\
      & \node [block] (combin) {\textsf{Solution Combination}};
      & \node [cloud] (subset) {\textsf{Subset Generation}};
      &
      \\
    };
    % connect all nodes defined above
    \begin{scope} [every path/.style=line]
      \path (init)     |-    (update);
      \path (update)   --    (newsol);
      \path (newsol)   --    node {no} (stop);
      \path (newsol)   --    node {yes} (subset);
      \path (subset)   --    (combin);
      \path (combin)   --    (improve);
      \path (improve)  --    (null1);
    \end{scope}
  \end{tikzpicture}
  \caption{Scatter Search General Scheme}
  \label{fig:SS_GS}
\end{figure}


\subsection{Diversity of a Solution}
% La diversidad de mide respecto al conjunto de referencia actual
The diversity of a solution,
is measured relative to the Reference Set;
% Se usa el perfect matching para calcular el costo
the perfect matching
is evaluated
with each one of the solutions in the Reference Set,
where the cost consist in
the sum of the distances
from each server of the given solution
and their match in the solution of the Reference Set.
% Se utiliza el minimo costo como medida de diversidad
The minimum cost perfect matching
is used as diversity value.
%\todo{Explicar como se calcula la diversidad}

\subsection{Scatter Search}
A Scatter-Search
with a dynamic reference set
was implemented,
and a path relinking method
is used as combination method.
We used a dynamic reference set
given that the number of solutions generated
in each combination is large.
To enhanced the solutions
two improvement methods are included,
one proposed by Berman \cite{berman1987stochastic},
and a proposed Local Search.
%\todo[inline]{Agregar background de Scatter Search}

In Algorithm \ref{alg:SSInitial} (see below)
a Reference Set is created,
consisting in the bests \textit{b} solutions,
and the most diverse \textit{b} solutions.

\begin{algorithm}
  \caption{Scatter Search Initial Phase} \label{alg:SSInitial}
  \begin{algorithmic}
    \Procedure{Initial Phase}{} 
    \State $Iteration \gets 0$
    \State $IterationsUnchanged \gets 0$
    \State $RefSet \gets EmptySet()$
    \Repeat
    \State $Sols \gets DiversificationGenerator()$
    \State $Improvement(Sols)$
    \Comment {either of the two methods proposed}
    \State $RefSet.Update(Sols)$
    \State $Iteration \gets Iteration + 1$
    \State $IterationsUnchanged \gets IterationsUnchanged + 1$
    \If {$HasNewSolutions(RefSet)$}
    \State $IterationsUnchanged \gets 0$
    \EndIf
    \Until {$IterationsUnchanged \geq MaxIterUnchanged$
      \textbf{or} $Iteration \geq MaxIter$}
    \EndProcedure
  \end{algorithmic}
\end{algorithm}


In Algorithm \ref{alg:SSPhase} (see below),
each pair of solutions
is combined using path relinking.
However,
since the Reference Set is dynamic
some solutions are displaced
before been combined.
\todo{Aclarar donde se manda llamar algoritmo 1}
\begin{algorithm}
  \caption{General Scheme of the Scatter Search} \label{alg:SSPhase}
  \begin{algorithmic}
    \Procedure{Scatter Search Phase}{$RefSet$}
    \Repeat
    \State $G_{X} \gets SubsetGeneration(RefSet)$
    \ForAll {$X \in G_X$}
    \If {$X \subset RefSet$}
    \Comment {Some element of $X$ may have been removed}
    \State $C_{X} = SolutionCombination(X)$
    \ForAll{$Sol \in C_X$}
    \State $Improvement(Sol)$
    \EndFor
    \State $RefSet.Update(C_{X})$
    \EndIf \EndFor
    \Until {$NoNewSolutions(RefSet)$}
    \EndProcedure
  \end{algorithmic}
\end{algorithm}


\section{Description of Components}
The components of Scatter Search template \cite{glover1998template}
consists of
specific subroutines of the following types:
\begin{itemize}
\item A Diversification Generator
\item An Improvement Method
\item A Reference Set Update Method
\item A Subset Generation Method
\end{itemize}
We described below
each of the components.

\subsection{A Diversification Generator}
To component is used to generate
a collection of diverse trial solutions,
using an arbitrary trial solution
(or seed solution) as an input.
%\begin{algorithm}
  \caption{Path Relinking Initial Phase}\label{pr_init}
  \begin{algorithmic}[0]
    \Procedure{PathRelinking}{$Instance$}
    \State $EliteSols \gets MultiStart(Instance)$
    \Repeat
    \State $MiscSols \gets GenerateMiscSols(Instance,EliteSols)$
    \ForAll{$x \in MiscSols$}
    \State $ImprovementMethod(x)$
    \State $EliteSols.Update(x)$
    \EndFor
    \Until{$PerfectMatchingCost(MiscSols,EliteSols) < \epsilon$}
    \State $SubsetControl(EliteSols)$
    \EndProcedure
  \end{algorithmic}
\end{algorithm}

In this research
we develop a Greedy Randomized Adaptive Search Procedure (GRASP)
as a diversification generator.
GRASP \cite{feo1995greedy}
is a multi-start procedure that combines
the power of greedy heuristics and randomization.

GRASP typically consists of two phases:
a greedy randomized construction phase
and an improvement phase.
In this case,
our proposed GRASP
consist only of the construction phase
because the improvement phase
is applied in other parts
of the scatter search algorithm.
\subsubsection{GRASP}
%Why did you choose GRASP?
We choose GRASP
because
in the proposed problem
any allocation of sites for adjusters
is a feasible solution,
so it was decided to give some intelligence
to this simple allocation
instead of choose
random points.
%What does it consist?
We start with a partial solution
(location of a smaller number of adjusters)
with only one adjuster allocated,
evaluate each site with the greedy function,
and choose one from the best $\alpha$ evaluations,
until each adjuster was locate.
%Which evaluation functions did you use?
Several functions were evaluated
for constructive algorithm
but because presented cooperativeness
it is difficult
to approximate the final results
in a partial solution
with a greedy function.
We use three greedy functions
to test the construction.
%Fist function
The first function,
was the p-mean
or the sum the distances
of each allocation 
from each located adjuster
with their nearest demand points.
\begin{equation}
  \label{eq:grasp1}
  \sum_{j=1}^{m}{
    \sum_{i=1}^{n}{
      t_{ij}y_{ij}^{1}
    }
  }
\end{equation}
%Second function
For the second function,
we try to incorporate cooperatively,
including in the evaluation
the distance of allocations
from each located adjuster,
with their $k$ nearest demand points
plus \textit{idle} probability
given that is allocated
to their $k-1$ nearest demand points.
\begin{equation}
  \label{eq:grasp2}
  \sum_{j=1}^{m}{
    \sum_{k=1}^{p}{
      \sum_{i=1}^{n}{
        h_{ij}^{k}t_{ij}y_{ij}^{k}
      }
    }
  }
\end{equation}
%Third function
The third function,
consist in use the Mean Service Time (MST) calibration method
proposed by Jarvis \cite{jarvis1985approximating},
to obtain more accuracy values
of the current mean response time.


\subsection{An Improvement Method}
The goal of this component
is to transform a trial solution
into one or more enhanced trial solutions. 
If no improvement
of the input trial solution results, 
the ``enhanced'' solution
is considered to be
the same as the input solution.

We developed two improvement methods,
the first one is an adaptation
of the method proposed by Berman et al. \cite{berman1987stochastic}
(see Appendix \ref{ch:Berman} for mode details)
and the second one,
a Local Search
described in the next section.
%\todo[inline]{Incluir descripcion del metodo de Berman}
\subsection{Proposed Local Search}
%the move
The move considered
is to relocate an agent from its current position
to a different position.
%the neighborhood
The entire neighborhood
consist of all possible moves
that can take place from a given solution.
However,
due to the high cost of move relocation assessment,
a reduced neighborhood is considered.
Thus,
instead of considering all possible moves,
the procedure focuses on relocating the adjuster
with the smallest workload
to a place around or near the adjuster with the largest workload.
\missingfigure[figwidth=6cm]{Example of less workload adjuster}
\missingfigure[figwidth=6cm]{Example of neighborhood to evaluate}


\subsection{A Reference Set Update Method}
This method is used
to build and maintain a Reference Set
consisting of the \textit{b} best solutions found
(where the value of \textit{b}
is typically small,
e.g., between 20 and 40),
organized in an array
to provide efficient accessing
by other parts of the method.

We choose
a two level Reference Set
to maintain
good quality solutions
and diverse solutions.
This is
because the combination method
generates similar solutions as the input,
%And we opt for be dynamic
%since the number
%of generated solutions
%is big (proportional to \textit{p}).
and Moreover,
the size of the reference set is dynamically determined,
according to $p$.

Whenever the Reference Set is update
with a set of solutions
%\todo{No se entiende!}
the update method
first try to incorporate the solutions
by quality criteria;
%because the Reference Set is divided
%in two parts,
%first try to accommodate each solution
if a solutions enters the Reference Set,
degrades a lower quality solution,
removing the worst solution of the Reference Set.
These steps are shown in Algorithm \ref{alg:rs_update}.
\begin{algorithm}
  \caption{Reference Set Update Method}\label{alg:rs_update}
  \begin{algorithmic}[0]
    \Procedure{ReferenceSet.Update}{$ToUpdateSols$}
    \State $Bin \gets PublicBin()$
    \For{$Sol \in ToUpdateSols$}
    \If{$Sol.Quality > RefSet.LowestQuality$}
    \State $DegradedSol \gets RefSet.LowestQualitySol()$
    \State $RefSet.Remove(DegragedSol)$
    \State $RefSet.InsertByQuality(Sol)$
    \State $Bin.Add(DegradedSol)$
    \Else
    \State $Bin.Add(Sol)$
    \EndIf \EndFor
    \EndProcedure
  \end{algorithmic}
\end{algorithm}


After adding solutions to the Reference Set
by using the quality criterion,
% Update Diversity
every solution that was generated
and did not enter the Reference Set,
displaced by a better solution,
or actually
a member of the diverse part of the Reference Set
is evaluated to be in the Reference Set
as a diverse solution.
As shown by the Algorithm \ref{alg:rs_updatediv}.
\begin{algorithm}
  \caption{Reference Set Update Diversity Method}
  \label{alg:rs_updatediv}
  \begin{algorithmic}[0]
    \Procedure{ReferenceSet.UpdateDiversity}{}
    \State $Bin \gets PublicBin()$
    \State $Bin \gets Bin \cup RefSet.DiversitySols()$
    \State $RefSet.Remove(RefSet.DiversitySols())$
    \For{$Sol \in Bin$}
    \State $Sol.DiversityValue \gets MinumumCostPerfectMatching(RefSet,Sol)$
    \EndFor
    \Repeat
    \State $DiverseSol \gets MostDiverse(Bin)$
    \State $Bin.Remove(DiverseSol)$
    \State $RefSet.InsertByDiversity(DiverseSol)$
    \For{$Sol \in Bin$}
    \State $Sol.UpdateDiversityValue(DiverseSol)$
    \EndFor
    \Until{$Count(RefSet.DiversitySols()) = b$ \textbf{or} $Count(Bin) = 0$}
    \EndProcedure
  \end{algorithmic}
\end{algorithm}

%solutions that do not enter the Reference Set
%and those that have been removed,
%are considered to be part
%of the Reference Set
%by using the diversity criterion.
%reassessing the diverse solutions,
%since they evaluation depends
%on the quality members.
%\todo[inline]{Agregar como se mide la diversidad}
%% We need layers to draw the block diagram
\pgfdeclarelayer{background}
\pgfdeclarelayer{foreground}
\pgfsetlayers{background,main,foreground}

% Define a few styles and constants
\tikzstyle{sensor}=[draw, fill=blue!20, text width=5em, 
  text centered, minimum height=2.5em]
\tikzstyle{ann} = [above, text width=5em]
\tikzstyle{naveqs} = [sensor, text width=6em, fill=red!20, 
  minimum height=12em, rounded corners]
\def\blockdist{2.3}
\def\edgedist{2.5}

\begin{figure}
  \label{fig:SS_Upd}
  \centering
  \begin{tikzpicture}
    \node (naveq) [naveqs] {Reference Set};
    % Note the use of \path instead of \node at ... below. 
    \path (naveq.140)+(-\blockdist,0) node (gyros) [sensor] {Quality};
    \path (naveq.-140)+(-\blockdist,0) node (accel) [sensor] {Diverse};
    
    %    \path [draw, ->] (gyros) -- node [above] {$\vc{\omega}_{ib}^b$} 
    %    (naveq.west |- gyros) ;
    % We could simply have written (gyros) .. (naveq.140). However, it's
    % best to avoid hard coding coordinates
    %      \path [draw, ->] (accel) -- node [above] {$\vc{f}^b$} 
    %      (naveq.west |- accel);
    \node (IMU) [below of=accel] {RefSet};
    \path (naveq.south west)+(-0.6,-0.4) node (INS) {INS};
    \draw [->] (naveq.50) -- node [ann] {Velocity } + (\edgedist,0) 
    node[right] {$\vc{v}^l$};
    \draw [->] (naveq.20) -- node [ann] {Attitude} + (\edgedist,0) 
    node[right] { $\mx{R}_l^b$};
    \draw [->] (naveq.-25) -- node [ann] {Horisontal position} + (\edgedist,0)
    node [right] {$\mx{R}_e^l$};
    \draw [->] (naveq.-50) -- node [ann] {Depth} + (\edgedist,0) 
    node[right] {$z$};
    
    % Now it's time to draw the colored IMU and INS rectangles.
    % To draw them behind the blocks we use pgf layers. This way we  
    % can use the above block coordinates to place the backgrounds   
    \begin{pgfonlayer}{background}
      % Compute a few helper coordinates
      \path (gyros.west |- naveq.north)+(-0.5,0.3) node (a) {};
      \path (INS.south -| naveq.east)+(+0.3,-0.2) node (b) {};
      \path[fill=yellow!20,rounded corners, draw=black!50, dashed]
      (a) rectangle (b);
      \path (gyros.north west)+(-0.2,0.2) node (a) {};
      \path (IMU.south -| gyros.east)+(+0.2,-0.2) node (b) {};
      \path[fill=blue!10,rounded corners, draw=black!50, dashed]
      (a) rectangle (b);
    \end{pgfonlayer}
  \end{tikzpicture}
  \caption{Scatter Search Update Method}
\end{figure}
\todo{Ajustar imagen si es que se va a dejar}


\subsection{A Subset Generation Method}
This component operates on the Reference Set,
%to produce a subset of its solutions
%as a basis for creating combined solutions.
it is used to create subsets
of solutions that are combined
by the Combination Method.
\begin{algorithm}
  \caption{Subsets Generator}\label{alg:subsets}
  \begin{algorithmic}[0]
    \Procedure{GenerateSubsets}{$RefSet,NowTime$}
    \State $NewSols \gets RefSet.SolutionsSince(NowTime)$
    \ForAll{$(Sol_x,Sol_y) \in NewSols$}
    \If{$Sol_x \in RefSet$ \textbf{and} $Sol_y \in RefSet$}
    \State $CombinedSols \gets PathRelinkingCombination(Sol_x,Sol_y)$
    \State $RefSet.Update(CombinedSols)$
    \EndIf \EndFor
    \State $RefSet.UpdateDiversity()$
    \If{$RefSet.NumberOfOldSols(NowTime) > 0$}
    \State $OldSols \gets RefSet.SolutionsUntil(NowTime)$
    \ForAll{$Sol_x \in NewSols$}
    \ForAll{$Sol_y \in OldSols$}
    \If{$Sol_x \in RefSet$ \textbf{and} $Sol_y \in RefSet$}
    \State $CombinedSols \gets PathRelinkingCombination(Sol_x,Sol_y)$
    \State $RefSet.Update(CombinedSols)$
    \EndIf \EndFor
    \State $RefSet.UpdateDiversity()$
    \EndFor \EndIf
    \EndProcedure
  \end{algorithmic}
\end{algorithm}

We only look for subsets of size two,
i.e. pairs of solutions;
because our Reference Set is dynamic,
at the beginning of each iteration
we label the solutions as new and old
to identify them,
and avoid create pairs previous combined.
To generate the pairs
we use
a new solution with another new solution,
%\todo{Refrasear}
after
if there are old solutions
we make pairs
with a new solution
and an old solution.

Is important to notice that
the new solutions that enters in the Reference Set
at the current iteration
are not part of the subsets
until the next iteration.

\subsection{A Solution Combination Method}
With this component,
our scatter search
transforms a given subset of solutions 
produced by the Subset Generation Method
into
one or more combined solution vectors.

We choose a path relinking as a combination method
consisting of
determining a match between servers,
this match can be
\begin{itemize}
\item Perfect Matching:
  minimizing the distance between paired servers
\item Workload Matching:
  sorting the servers of both solutions
  according to their workload,
  and match them according to these sorted lists.
\item Random Matching
\end{itemize}
Once we have the matching
we proceed from one solution
to interchange each server,
one by one
until the other solution is reached;
in each interchange
we have a new solution.
To make the interchanges
we have three options:
\begin{itemize}
\item Nearest First
\item Farthest First
\item Random
\end{itemize}
\begin{algorithm}
  \caption{Path Relinking Combination Method}\label{alg:pr_combination}
  \begin{algorithmic}[0]
    \Procedure{PahtRelinkingCombination}{$Sol_x,Sol_y$}
    \State $CombinedSols \gets EmptyList()$
    \State $match \gets Matching(Sol_x,Sol_y)$
    \Comment{perfect,workload,random}
    \State $order \gets ProcessOrder(Sol_x,match,Sol_y)$
    \Comment{nearest/farthest first,random}
    \For{$i \gets 1,p$}
    \State $j \gets order[i]$
    \If{$Sol_x.ServerLocation(j) != Sol_y.ServerLocation(match[j])$}
    \State $Sol_x.SetServerLocation(j,Sol_y.ServerLocation(match[j]))$
    \State $CombinedSols.insert(Sol_x)$
    \EndIf \EndFor
%    \State \textbf{return} $CombinedSols$
    \EndProcedure
  \end{algorithmic}
\end{algorithm}

%\todo[inline]{Hacer mas claro la secuencia y metodos}


\chapter{Evaluation}
\section{Model A}
%This model was tested 
%with a benchmark instance of 30 points, 
%which sets $V$ and $W$ are the same.
%For values of $p = 1,\ldots,7$
%the optimal solution was found in less than one minute.
%
%For other benchmark instances
%was impossible to find the solution
%because they have more than 300 points,
%and the proposed model grows to over one million variables.
%  For some large instances ($n = 50$ , $m = 30$)%

\todo[inline]{Computational Results}

\section{Model B}

\section{GRASP}
Several instances were tested
with different $\alpha$ values
to determine the value of $\alpha$ to use.
The value of $\alpha = 0$ means that
the procedure becomes completely greedy,
in contravention
a value of $\alpha = 1$ means that
the procedure becomes completely random.
\missingfigure[figwidth=6cm]{GRASP results of Objective versus $\alpha$}

The test suggest
that
with a more random version
there are more different solutions
and a wide range of objective values
achieving better solutions.

\section{Scatter Search}
To evaluate this heuristic
several instances was tested
and
a desing of experiment was made
to determine
the match method
and the processing order method
that produce better results.

\begin{center}
  \begin{tabular}{|l|l|l|r|r|r|}
    \hline
    M & N & p & Improvement & Time (sec) & MST(sec) \\ \hline
    50 & 30 & 7 & 6.97 & 1.55 & 1.64 \\
    &  & 10 & 9.62 & 5.20 & 4.57 \\
    &  & 15 & 6.95 & 8.55 & 6.49 \\
    &  & 20 & 3.67 & 13.79 & 8.80 \\ \cline{2-6}
    & 50 & 7 & 9.76 & 2.47 & 2.41 \\
    &  & 10 & 12.20 & 7.09 & 6.15 \\
    &  & 15 & 9.75 & 11.35 & 8.51 \\
    &  & 20 & 10.63 & 30.22 & 19.38 \\ \cline{2-6}
    & 75 & 7 & 11.64 & 2.43 & 2.37 \\
    &  & 10 & 14.50 & 6.16 & 5.35 \\
    &  & 15 & 18.26 & 19.44 & 14.08 \\
    &  & 20 & 13.59 & 44.19 & 28.05 \\ \hline
    Total Result &  &  & 10.63 & 12.71 & 8.98 \\ 
    \hline
  \end{tabular}
\end{center}

\begin{center}
  \begin{tabular}{|l|l|l|r|r|r|}
    \hline
    M & N & p & Improvement & Time (sec) & MST(sec) \\ \hline
    100 & 30 & 7 & 6.70 & 4.10 & 4.33 \\
    &  & 10 & 5.74 & 9.97 & 9.65 \\
    &  & 15 & 3.38 & 14.90 & 13.46 \\
    &  & 20 & 2.80 & 28.30 & 24.01 \\ \cline{2-6}
    & 50 & 7 & 7.68 & 4.23 & 4.46 \\
    &  & 10 & 8.09 & 11.01 & 10.59 \\
    &  & 15 & 9.55 & 25.72 & 22.36 \\
    &  & 20 & 9.51 & 61.22 & 49.41 \\ \cline{2-6}
    & 75 & 7 & 7.70 & 4.55 & 4.76 \\
    &  & 10 & 11.46 & 10.17 & 9.82 \\
    &  & 15 & 11.95 & 36.70 & 31.19 \\ \hline
    Total Result &  &  & 7.69 & 19.17 & 16.73 \\ 
    \hline
  \end{tabular}
\end{center}

\missingfigure[figwidth=6cm]{Scatter Search results}


\chapter{Conclusions}
\begin{itemize}
\item Two different models (Model A and Model B)
  for the car wreck adjusters location have been proposed. 
  The models can be seen as linear approximations to the SQpM.
\item We observed that the Model A
  presents more difficulties for solving medium instances.
\item For small enough values of $\ell$,
  Model B can be solved relatively quick
  for instances of medium size.
\item The solutions obtained from the Model B
  for small values of $\ell$
  approximates well the solutions of Model A,
  as well as self solutions for large values of $\ell$
\end{itemize}


\appendix
%%% Haz un documento para cada apéndice
% Approximating the Equilibrium Behavior of Multi-Server Loss Systems
% Jarvis - Management Sience
% 1985

\section{Approximating the Equilibrium Behavior of MSLS}
In Jarvis \cite{jarvis1985approximating}
a procedure is given
for approximating the equilibrium behavior
of multi-server loss systems having distinguishable servers
and multiple customers types
under light to moderate traffic intensity.

\subsection{Introduction}
In an emergency service such as fire or police, 
the servers are fire fighting units or patrol cars
and the customers are calls for service.
The simple Erlang loss system
is inadequate in two aspects for a detailed system analysis.
\begin{enumerate}
\item one often wishes
  to preserve the identity of service units (distinguishable servers).
\item because
  of the geographic nature of these systems,
  the service time depend on both
  the server and the customer
  at least through the travel time between the pair.
\end{enumerate}

\subsection{Model assumptions, Notation, and Terminology}
Consider a system in which:
\begin{enumerate}
\item Exactly one server
  is assigned to each customer
  unless all servers are busy, 
  in which case
  the customer is irrevocably lost of the system
\item Servers are assigned
  to customers according to a fixed preference assignment rule
\item No preemption of service is allowed
\item Assignments are made immediately
  upon customer arrival
\end{enumerate}

and we have the following parameters
\begin{enumerate}
\item \textit{N} distinguishable servers
\item \textit{C} types of customers
\item Customers of type \textit{m}
  arrive according to a Poisson process with rate $\lambda_{m}$
\item $\lambda$ total arrival rate
\item $a_{mk}$ be the \textit{k}th preferred server
  for customers of type \textit{m}
\item $\tau_{im}$ the expected service time
  for server \textit{i} and customer of type \textit{m}
\end{enumerate}

The performance measures for the system include 
\begin{enumerate}
\item $\rho_i$ the workload of server \texit{i}
\item $f_{im}$ the probability a random customer of type \textit{m}
  is assigned to server \textit{i}
\item $P_{N}$ the probability all servers are busy
\end{enumerate}

\subsection{Approximation Procedure}
The procedure described is based on that given by
Larson \cite{larson1975approximating},
for approximating performance measures
for the Hypercube model
assuming exponential service times.
Larson developed an approximation for $f_{im}$ as
\begin{equation} \label{f_im}
  f_{im} \simeq Q(N,p,k-1)(1-\rho_{i})\prod_{l=1}^{k-1}{\rho_{a_{ml}}}
\end{equation}
where
\begin{equation} \label{Q}
  Q(N,p,k) =
  \sum_{j=k}^{N-1}{
    \frac{(N-j)(N^j)(\rho^{j-k})P_0(N-k-1)!}
         {(j-k)!(1-P_N)^kN!(1-\rho(1-P_N))}
  } \\
  \hfill \mbox{for} k = 0,1,\ldots,N-1
\end{equation}

Let $B_i$ denote the event
that server \textit{i} is busy;
and let $B_{im}$ denote the event
that server \textit{i} is busy
serving a customer of type \textit{m}, then
\begin{equation} \label{rho_i}
  \rho_{i} = 
  Pr\left[B_i\right] = 
  \sum_{m=1}^{C}{
    Pr\left[B_{im}\right]} =
  \sum_{m=1}^{C}{
    \lambda_{m}f_{im}\tau_{im}
  }
\end{equation}
combine equations (\ref{f_im}) and (\ref{rho_i})
and solve for $\rho_i$ to obtain the approximation iteration
\begin{equation}
  \rho_i\mbox{(new)} =
  \frac{V_i}
       {(1+V_i)}
\end{equation}
where $V_i$ is given by
\begin{equation} \label{V_i}
  V_i = \sum_{k=1}^{N}{\sum_{m:a_{mk}=i}{\lambda_m \tau_{im} Q(N,\rho,k-1)\prod_{l=1}^{k-1}{\rho_{a_{ml}}}}}
\end{equation}

When there is a common mean service time, the estimates for $\rho_i$ can be normalized using
\begin{equation} \label{P_N}
  \sum_{i=1}^{N}{\rho_i} = N \rho (1 - P_{N})
\end{equation}

In the generalized procedure, $\tau$ can be approximated at the end of each iteration by
\begin{equation} \label{tau}
  \tau = \sum_{m=1}^{C}{\left(\frac{\lambda_m}{\lambda}\right)\sum_{i=1}^{N}{\frac{\tau_{im}f_{im}}{(1-P_N)}}}
\end{equation}

{\footnotesize
  \begin{center}
    \textit{ Approximation Algorithm}
  \end{center}
  \vspace{-8pt}
  \hline \\
  \vspace{2pt}
  \textit{Given:}
  \begin{equation*}
    \lambda_m,\tau_{im},a_{mk}\mbox{  for } m = 1,\ldots,C;i=1,\ldots,N;k=1,\ldots,N \hfill
  \end{equation*}
  \textit{Initialize:}
  \begin{equation*}
    \rho_i = \sum_{m: a_{m1} = i}{\lambda_m \tau_{im}}; \; \tau = \sum_{m=1}^{C}{(\lambda_m/\lambda)\tau_{a_{m1},m}} \hfill
  \end{equation*}
  \textit{Iteration:}\\
  (1) Compute $Q(N,\rho,k)$ for $k = 1,\ldots,N-1$ where $\rho = \lambda \tau / N$ using equation (\ref{Q}). \\
  (2) For $i = 1,\ldots,N$, the new $\rho_i$ is $V_i/(1+V_i)$, where $V_i$ is given by equation (\ref{V_i}). \\
  (3) Stop if max change in $\rho_i$ is less than convergence criterion. \\
  (4) Else compute $P_N$ by equation (\ref{P_N}), $\tau$ by equation (\ref{tau}), and $f_{im}$ by equation (\ref{f_im}). \\
  (5) Return to step 1. \\
  \hline
}

No analytic bounds on the accuracy or convergence properties of the approximation procedure have been developed.

In regards to convergence properties,
the numerical iteration has proved to be very stable and converges in a small number of iterations under relatively stringent conditions,
with 4 to 6 iterations being typical for 10 servers systems.

In comparing the accuracy of this approximation to results of the exact Hypercube model,
Larson has found errors in server workloads to be less than 1 to 2 percent.


\backmatter
\pagestyle{main}

%%% Aquí va la bibliografía, puedes usar el entorno de LaTeX (thebibliography)
%%% o la herramienta BibTeX. En caso de que optes por BibTeX, puedes usar
%%% alguno de los archivos de estilo (mighelbib.bst o mighelnat.bst) incluidos
%%% en el paquete, cuyos diseños armonizan con el diseño de tesis provisto por
%%% fime.cls. Para muestra, basta un botón:
\bibliographystyle{abbrv}
\bibliography{references}

\label{lastpage}
%Autobiografia

\chapter*{Resumen autobiográfico}
\thispagestyle{empty}

\begin{center}
\autor

Candidato para obtener el grado de\\
\grado\\
\orientacion\bigskip

\uanl\\
\fime\bigskip

Tesis:\\
\textsc{\large\titulo}
\end{center}\bigskip

%Aquí va tu historia
Aquí va tu historia. Recuerda que debe incluir: lugar y fecha de nacimiento, nombre de los padres, escuelas y universidades en las que se graduó después de la preparatoria, títulos o grados obtenidos (no incluir los estudios que se están concluyendo), experiencia profesional y organizaciones profesionales a las que pertenece (no incluir lista de publicaciones).


\end{document}
